\chapter{Сводка алгоритмов}
\label{ch:algorithms}

Псевдокод ключевых операторов, собранный в одном месте для удобства. Это скелеты,
а не финальные реализации: они фиксируют логику, обсуждённую в главах, но опускают
инженерные детали (батчинг мутаций, кэширование $\pi$, индексы эмбеддингов).

\section{Единый блуждающий агент (три политики)}

\begin{lstlisting}[style=pseudo]
def walk(G, start_nodes, mode, steps, mood, query=None):
    # mode задаёт политику через параметры: bias перехода, alpha (телепорт),
    # bias_uncertainty. Один движок — три режима (глава про обход).
    if mode == "retrieval":   alpha, bu = 0.0, 0.0
    if mode == "curiosity":   alpha, bu = 0.05, 1.0     # тянет к need_to_discover
    if mode == "sleep":       alpha, bu = 0.30, 0.0     # телепорт-скраб

    active = set(start_nodes); visited = Counter()
    for _ in range(steps):
        A = sample(active)
        if random() < alpha:                 # телепорт (рестарт)
            B = random_node(G)
        else:
            # $P(A\to B)\propto w_{AB}\cdot rel(B|query)\cdot moodcong(B|mood)\cdot unc(B)^{bu}$
            B = sample_neighbor(G, A, weights=transition_scores(
                    G, A, query, mood, bias_uncertainty=bu))
        visited[B] += 1
        maybe_fire_operations(G, B, mode)    # retrieval / fusion / inspect
        active = frontier(active, B)
    return visited
\end{lstlisting}

\section{Локальное правило и лавинная пропагация}

\begin{lstlisting}[style=pseudo]
def propagate(G, A, sigma_target=0.9):
    # Пропагация нового факта A как затухающей лавины с контролем ветвления.
    frontier = link_initial(G, A)            # A -> N(A), напр. последние top-k
    while frontier:
        next_frontier = []
        # нормировка под целевой ветвящийся коэффициент sigma < 1
        gain = branching_normalizer(frontier, sigma_target)
        for B in frontier:
            d = delta_conf(G, A, B) * gain    # см. ниже: знак, источник, независимость
            G.confidence[B] = clip(G.confidence[B] + d, 0, 1)
            th_e = theta_emergent(G, B)       # растёт с глубиной консолидации
            if abs(d) > th_e:                 # крупная переоценка -> в ширину
                next_frontier += neighbors(G, B)
            elif abs(d) > THETA_TEMP:         # промежуточная зона -> пометить
                mark(G, B, "need_to_revise")
            # else: глохнет здесь
        frontier = next_frontier

def delta_conf(G, A, B):
    sign      = edge_sign(G, A, B)            # supports (+) / contradicts (-)
    strength  = edge_weight(G, A, B)
    gate      = G.confidence[A] * source_trust(G, A)   # карантин нового факта
    indep     = independence_discount(G, A, B)         # анти-эхо по провенансу
    return ETA * sign * strength * gate * indep * G.confidence[A]
\end{lstlisting}

\section{Забывание, прунинг, консолидация}

\begin{lstlisting}[style=pseudo]
def nightly_maintenance(G):
    for v in G.nodes:
        G.strength[v] = act_r_base_level(G.access_log[v])   # частота + свежесть
    for v in low_activation(G):
        if importance(G, v) low and confidence(G, v) low:
            drop(G, v)                                       # редкое-и-неважное
    for cluster in redundant_episodic_clusters(G):
        merge_to_semantic(G, cluster)   # огрубление: эпизоды -> обобщение
    # эмбриональным узлам — меньшая скорость распада (континуум глубины)
\end{lstlisting}

\section{Телепорт-скраб (анти-энтропия)}

\begin{lstlisting}[style=pseudo]
def scrub(G, budget):
    # Случайная инспекция: гарантия покрытия за счёт эргодичности телепорта.
    for _ in range(budget):
        v = random_node(G)                    # рестарт куда угодно
        for w in local_walk(G, v, depth=k):
            if contradicts_any(G, w):         # столкновение -> ревизия
                demote(G, w); mark(G, justificands(G, w), "need_to_revise")
            if flagged(G, w, "need_to_discover"):
                maybe_fetch_external(G, w)     # любопытство/внешний вход
\end{lstlisting}

\section{Настроение и монитор руминации}

\begin{lstlisting}[style=pseudo]
def mood_update(mood_layers, event):
    # Стек текущих интеграторов с разными tau; актуальное = суперпозиция.
    for L in mood_layers:                     # session, day, ..., life
        mood_L = (1 - 1/L.tau) * mood_L + (1/L.tau) * appraise(event)
    mood_now = sum(w_L * mood_L for L in mood_layers)
    if autocorr(history(mood_now)) > RUMINATION_TH:   # залипла валентность
        trigger(one_of=["reframe", "attention_shift", "invoke_sleep"])
    return mood_now
\end{lstlisting}

\section{Оценивание, значимость и консолидация реакции}

\begin{lstlisting}[style=pseudo]
def encode_event(G, event, mood_prev, mood_now):
    ap = appraisal(event)                     # surprise, goal_congruence, valence...
    salience = ap.surprise * ap.goal_relevance
    #  дешёвый прокси значимости, почти бесплатный:
    salience_cheap = abs(mood_now - mood_prev)
    consolidate(G, event, weight=salience)    # гейт значимости, НЕ одобрения
    if ap.surprise high: mark(G, event, "need_to_discover")

def update_reaction(G, situation, user, reaction, received_ok):
    e = reaction_edge(G, situation, user, reaction)   # условь по (ситуация, отношение)
    if received_ok: e.beta.alpha += 1
    else:           e.beta.beta  += 1                 # угасание, не второй коммит
    e.confidence = e.beta.alpha / (e.beta.alpha + e.beta.beta)
    if e.confidence > CONSOLIDATE_TH: promote_to_slow_layer(G, e)
\end{lstlisting}

\section{Скорер характера и отбор голоса}

\begin{lstlisting}[style=pseudo]
def phi(text, axes):                          # LLM-судья по рубрикам осей
    return [judge(text, rubric_a) for rubric_a in axes]   # вектор в R^d

def select_voice_corpus(candidates, target_t, axes):
    scored = [(c, phi(c, axes)) for c in candidates]
    # отбор/перевзвешивание, чтобы среднее выборки приблизилось к target_t
    return greedy_match_mean(scored, target_t)   # не усреднять муз, а фильтровать
\end{lstlisting}

\section{Стенд эксперимента}

\begin{lstlisting}[style=pseudo]
def run_experiment(build_graph, plant, operator, metric):
    G = build_graph()          # ярус 1: SBM/BA/WS + подсаженные сценарии
    truth = plant(G)           # ground truth: что должно произойти
    operator(G)                # walker / propagate / scrub / decay ...
    return metric(G, truth)    # precision/recall починки, сходимость, покрытие
\end{lstlisting}
